\documentclass[12pt,a4paper]{article}
\usepackage[utf8]{inputenc}
\usepackage{amsmath,amsfonts,amssymb}
\usepackage{graphicx}
\usepackage{booktabs}
\usepackage{hyperref}
\usepackage{geometry}
\geometry{margin=2.5cm}
\usepackage{caption}
\usepackage{float}
\usepackage{enumitem}

\title{\textbf{EmotionLens: Cross-Domain Facial Expression Recognition\\
with Controlled-Variable Experiments}}
\author{Group X \\ XMUM 2604 Natural Language Processing / Deep Learning}
\date{June 2026}

\begin{document}

\maketitle

\begin{abstract}
We present EmotionLens, a cross-domain facial expression recognition (FER) system built on a ResNet18 backbone trained on three complementary data sources: FERPlus with 10-voter soft labels, RAF-DB with real-world faces, and a self-captured dataset of 1,123 manually reviewed frames from six contributors. Through six controlled experiments---each altering exactly one variable---we isolate and quantify the contributions of architecture (VGG16$\to$ResNet18: $-$2.3pp), label quality (hard labels$\to$10-voter soft labels: +35.2pp), and data diversity (single-source$\to$multi-source: +20.0pp) to real-face generalization. Our key finding is that label quality is the single largest lever for cross-domain FER performance, followed by data diversity; architecture alone contributes essentially nothing. The best model achieves 89.0\% accuracy on the FERPlus test set and 64.5\% on unseen real faces---a 52.9pp improvement over the FER2013+VGG16 course baseline (11.6\%). The model is deployed as a real-time web application with five swappable analysis lenses, embodying a design philosophy in which the inference engine is the reusable core and application logic is open-ended.
\end{abstract}

\tableofcontents
\newpage

%=====================================================================
\section{Introduction}
%=====================================================================

Facial expression recognition (FER) aims to automatically classify human emotions from facial images. Based on Ekman and Friesen's (1971) seven basic emotion categories---neutral, happiness, surprise, sadness, anger, disgust, and fear---FER enables a broad range of applications, from classroom engagement monitoring and human-computer interaction to online proctoring and mental health screening.

While modern deep networks routinely exceed 85\% accuracy on standard FER benchmarks (Barsoum et al., 2016; Li \& Deng, 2019), a critical bottleneck persists: models trained on curated laboratory datasets often collapse when deployed on real faces they have never seen. This phenomenon---variously termed domain gap or subject identity bias---means the model may learn to recognize who a person is rather than what emotion they are expressing.

In this course, the instructor suggested starting from the FER2013 dataset with a VGG16 backbone. We faithfully reproduced this baseline: 69.5\% on the FER2013 standard test set, but only 11.6\% on our own self-captured dataset of 1,123 real-face images---barely above random chance. This 58-percentage-point performance gap became the starting point of our investigation.

Our work addresses three questions: (1) What factors determine a model's ability to generalize to real, unseen faces? (2) Can we isolate and quantify each factor's contribution through strict controlled-variable experiments? (3) What does a model that works on real faces unlock in terms of deployable applications?

To answer these questions, we designed six experiments that form a causal chain: each step changes exactly one variable---architecture, label quality, or data diversity---while holding all others constant. The resulting decomposition reveals where the gains actually come from. We then deploy the best model as EmotionLens, a real-time web application with five swappable analysis lenses that demonstrate the engine's generality.

The remainder of this paper is organized as follows. Section 2 reviews related work. Section 3 describes our datasets. Section 4 details the model implementation and experiment design. Section 5 presents results and analysis. Section 6 concludes.

%=====================================================================
\section{Literature Review}
%=====================================================================

\subsection{Evolution of FER Datasets}

FER2013 (Goodfellow et al., 2013) was the first large-scale FER benchmark, comprising 35,887 $48\times48$ grayscale images, each assigned a single hard label (0--6) by one annotator. However, single-annotator subjectivity and the inherent ambiguity of facial expressions lead to substantial label noise in FER2013.

FERPlus (Barsoum et al., 2016) reuses the exact same image set but changes the annotation protocol: each image is independently labeled by 10 crowd-workers, producing a probability distribution over emotion classes (e.g., $[0.0, 0.1, 0.7, 0.1, 0.0, 0.0, 0.1]$). These soft labels significantly reduce annotation noise and add a contempt class, forming an 8-class taxonomy. In our work, we discard contempt (which has insufficient samples in RAF-DB and our self-captured set) and unify all sources to seven Ekman classes.

RAF-DB (Li \& Deng, 2019) is a real-world FER dataset with 15,344 well-aligned, high-resolution face images. Unlike FERPlus, whose images are predominantly low-resolution web crops, RAF-DB covers diverse real-world conditions---varying poses, illuminations, and backgrounds. In our experiments, RAF-DB serves as a domain bridge that compensates for FERPlus's lack of real-world diversity.

\subsection{Architecture Choices}

VGG16 (Simonyan \& Zisserman, 2014) is a classic CNN with 138M parameters and a simple stacked-convolution design. While historically important, its large parameter count and slow inference make it suboptimal for real-time applications.

ResNet (He et al., 2016) introduced residual connections that mitigate vanishing gradients, enabling stable training of deeper networks. ResNet18, with only 11.2M parameters (roughly 1/12 of VGG16), benefits from ImageNet pre-training and transfers well to downstream tasks.

EfficientNet (Tan \& Le, 2019) uses neural architecture search to balance depth, width, and resolution. Despite its theoretical efficiency, our preliminary experiments with EfficientNet-B0 under the same training recipe yielded only 41.0\% on the FERPlus test set---far below ResNet18's 79.9\%. We therefore do not use EfficientNet as our primary architecture.

Our architecture choice is empirical rather than dogmatic: we benchmarked VGG16, ResNet18, and EfficientNet-B0 under identical training conditions. ResNet18 offers the best practical trade-off among parameter efficiency, training stability, and final performance.

\subsection{Cross-Domain Generalization and Regularization}

Domain shift---the mismatch between training and deployment distributions---is the central challenge in deploying FER systems. A model trained on $48\times48$ web-crawled grayscale images faces a severe distribution shift when applied to high-resolution webcam captures in indoor lighting. Subject identity bias compounds this problem: when a model sees multiple frames of the same person during training, it may learn identity-specific features rather than emotion-generic features (Chen et al., 2022).

We employ several regularization techniques to mitigate these issues. Mixup (Zhang et al., 2018) trains on convex combinations of image pairs and their labels, improving robustness to distribution shift. Label smoothing (Szegedy et al., 2016) softens the training targets, preventing the model from becoming overconfident on noisy labels. AdamW (Loshchilov \& Hutter, 2019) decouples weight decay from adaptive gradient scaling, providing better generalization than standard Adam.

%=====================================================================
\section{Data Set}
%=====================================================================

\subsection{Public Datasets}

Our training corpus draws from three public sources:

\textbf{FERPlus} (Barsoum et al., 2016): The official Microsoft release in parquet format. 35,803 images (train 28,709 / validation 3,546 / test 3,546), originally 8-class soft labels. We drop the contempt class to unify with our other sources, resulting in 7-class probability distributions. Training uses soft cross-entropy loss against the 10-voter distribution.

\textbf{RAF-DB} (Li \& Deng, 2019): Approximately 15,344 real-world face images with 7-class hard labels (basic set only). RAF-DB faces are better aligned and cover more diverse real-world conditions than FERPlus.

\textbf{FER2013} (Goodfellow et al., 2013): 35,887 $48\times48$ grayscale images in CSV format with 7-class hard labels. We use FER2013 only in Exp1 (VGG16 baseline) and Exp1b (ResNet18 architecture ablation) to maintain consistency with the course baseline.

\subsection{Self-Captured Data}

To evaluate generalization to real faces---specifically, faces of our own group members---we collected a self-captured dataset.

\textbf{Collection protocol}: Six group members each recorded short video sessions on their personal laptop webcams under ordinary indoor conditions (classrooms or dormitories, natural lighting, seated approximately 40--60 cm from the screen). For each of the seven target classes, the contributor was asked to pose the corresponding expression while the recording script saved frames at a fixed interval.

\textbf{Annotation workflow}: After recording, frames were extracted at regular intervals. A small pre-trained FER model produced weak labels as an initial guess. Group members then manually reviewed every frame---correcting errors, removing blurred or ambiguous frames. The final reviewed set contains 1,123 frames.

\textbf{Split}: We use an 80/20 stratified random split (train/validation). We acknowledge that a subject-disjoint split (where the same person never appears in both train and test) would be methodologically preferable, and we discuss this in Section 6 as a limitation.

\subsection{Preprocessing}

All images undergo a unified preprocessing pipeline: (1) resize to $224\times224$ RGB; (2) normalize with ImageNet statistics ($\mu=[0.485, 0.456, 0.406]$, $\sigma=[0.229, 0.224, 0.225]$). For FER2013/FERPlus $48\times48$ grayscale originals, we resize and replicate to three channels. RAF-DB RGB images are directly resized. Self-captured frames are first passed through MTCNN (Zhang et al., 2016) for face detection and cropping, then resized.

\begin{table}[H]
\centering
\caption{Dataset statistics}
\begin{tabular}{lcccc}
\toprule
Source & Train & Validation & Test & Total \\
\midrule
FER2013 & 28,709 & 3,589 & 3,589 & 35,887 \\
FERPlus & 28,709 & 3,546 & 3,546 & 35,803 \\
RAF-DB  & $\sim$12,000 & — & $\sim$3,000 & $\sim$15,344 \\
Self    & 899 & 224 & — & 1,123 \\
\bottomrule
\end{tabular}
\end{table}

%=====================================================================
\section{Implementation}
%=====================================================================

\subsection{Model Architecture}

We adopt ResNet18 (He et al., 2016) as our primary backbone, initialized with ImageNet pre-trained weights from torchvision. The original 1000-class fully connected layer is replaced with Linear(512, 7), mapping the 512-dimensional pooled feature vector to seven emotion logits. Inputs are $224\times224$ RGB images with ImageNet normalization. The model has 11,180,103 trainable parameters.

The choice of ResNet18 over VGG16 and EfficientNet-B0 is empirically grounded: under identical training recipes, VGG16 (138M parameters) achieves 69.5\% on FER2013 with slower inference; EfficientNet-B0 (5.3M parameters) reaches only 41.0\% on FERPlus. ResNet18 offers the best practical trade-off in parameter efficiency, training stability, and final performance. The residual connections also mitigate vanishing gradients, enabling stable convergence within 35 epochs.

\subsection{Training Recipe}

\subsubsection{Loss Function}

We use two loss functions depending on the label type of the training data:

\textbf{Hard-label Cross-Entropy} (FER2013 and RAF-DB):
\begin{equation}
\mathcal{L}_{\text{hard}} = -\log p_{\text{model}}(y_{\text{true}})
\end{equation}
This loss forces the model toward 100\% confidence in a single label, which can be problematic when annotation noise is present.

\textbf{Soft Cross-Entropy} (FERPlus):
\begin{equation}
\mathcal{L}_{\text{soft}} = -\sum_{k=1}^{K} p_{\text{voter}}(k) \cdot \log p_{\text{model}}(k)
\end{equation}
where $p_{\text{voter}}(k)$ is the fraction of 10 annotators who voted for class $k$. This allows the model to match a distribution rather than a point estimate, making it more robust to annotation noise. The contrast between these two loss formulations is itself part of our experimental design---it directly quantifies the impact of label quality.

\subsubsection{Optimizer and Gradient Handling}

We use AdamW (Loshchilov \& Hutter, 2019) with the update rule:
\begin{equation}
\theta_{t+1} = \theta_t - \eta \cdot \left(\frac{\hat{m}_t}{\sqrt{\hat{v}_t} + \epsilon} + \lambda \theta_t\right)
\end{equation}
where $\hat{m}_t$ and $\hat{v}_t$ are bias-corrected first- and second-moment estimates, and $\lambda$ is the decoupled weight decay coefficient. Hyperparameters: $\beta_1=0.9$, $\beta_2=0.999$, $\eta=3\times10^{-4}$, $\lambda=1\times10^{-4}$. Compared to standard Adam, AdamW decouples weight decay from the adaptive learning rate, providing faster convergence than SGD with better generalization than vanilla Adam.

Training uses Automatic Mixed Precision (AMP) via \texttt{torch.amp.GradScaler}: forward passes run in FP16 for speed, the loss is scaled before backpropagation to prevent gradient underflow, and gradients are unscaled before the optimizer step.

\subsubsection{Learning Rate Schedule}

We use CosineAnnealingLR:
\begin{equation}
\eta_t = \eta_{\min} + \frac{1}{2}(\eta_{\max} - \eta_{\min})\left(1 + \cos\left(\frac{t}{T_{\max}}\pi\right)\right)
\end{equation}
with $T_{\max}=35$, $\eta_{\max}=3\times10^{-4}$, $\eta_{\min}\approx 0$. Cosine annealing requires no manual step tuning and provides a smooth decay that allows later epochs to explore finer loss landscapes.

\subsubsection{Sampling}

To address class imbalance (e.g., happiness far outnumbers disgust), we use WeightedRandomSampler with per-class weights $w_c = 1 / \text{count}(c)$, replacement enabled, sampling the full dataset size each epoch. This oversamples minority classes (disgust, fear) and undersamples majority classes (happiness, neutral), ensuring balanced gradient signals.

\subsubsection{Regularization}

\textbf{Label Smoothing} ($\varepsilon=0.1$): Softens hard labels to $y_{\text{smooth}} = (1-\varepsilon) \cdot y_{\text{hard}} + \varepsilon/K$, preventing overconfidence.

\textbf{Mixup} ($\alpha=0.2$): Trains on convex combinations $\tilde{x} = \lambda x_i + (1-\lambda) x_j$, $\tilde{y} = \lambda y_i + (1-\lambda) y_j$ with $\lambda \sim \text{Beta}(\alpha,\alpha)$. Mixup improves robustness to distribution shift (Zhang et al., 2018).

\textbf{Weight Decay} ($1\times10^{-4}$): Applied through AdamW's decoupled mechanism, constraining weight magnitude.

\subsubsection{Data Augmentation}

Training augmentations: RandomHorizontalFlip ($p=0.5$), ColorJitter (brightness/contrast/saturation each 0.2), RandomErasing ($p=0.1$, scale=(0.02, 0.1)). Test-time augmentation uses horizontal-flip averaging.

\begin{table}[H]
\centering
\caption{Unified training hyperparameters across all experiments}
\begin{tabular}{lll}
\toprule
Hyperparameter & Value & Rationale \\
\midrule
Optimizer & AdamW ($\beta_1$=0.9, $\beta_2$=0.999) & Decoupled weight decay \\
Learning rate & $3\times10^{-4}$ & Empirically faster convergence than $1\times10^{-4}$ \\
Weight decay & $1\times10^{-4}$ & Standard value \\
LR schedule & CosineAnnealing, $T_{\max}$=35 & Smooth, no manual tuning \\
Epochs & 35 & Validation accuracy saturates near this point \\
Batch size & 128 (VGG16: 64) & VGG16 limited by 8GB VRAM \\
Loss & CE (hard) / Soft CE & Matches label type \\
Label smoothing & 0.1 & Inception-v3 recommendation \\
Mixup $\alpha$ & 0.2 & Conservative mixing \\
Hardware & RTX 4060 Laptop (8GB) & Local training \\
\bottomrule
\end{tabular}
\end{table}

\subsection{Experiment Design}

The experiment design is the methodological core of this work. We designed six experiments in which each step changes \textbf{exactly one variable}, enabling clean causal attribution of every performance change.

\begin{table}[H]
\centering
\caption{Experiment design: controlled-variable matrix}
\begin{tabular}{lccccl}
\toprule
Exp & Architecture & Training Data & Label Type & Variable Changed & Question Answered \\
\midrule
1  & VGG16    & FER2013    & Hard & --- (baseline)     & Course baseline performance? \\
1b & ResNet18 & FER2013    & Hard & Architecture       & Architecture contribution? \\
2  & ResNet18 & FERPlus    & Soft & Label quality      & Soft-label improvement? \\
3  & ResNet18 & F+R+S      & Mixed& Data diversity    & Added-data improvement? \\
4  & All 4    & ---        & ---  & Cross-domain eval  & Generalization patterns? \\
5  & FER lib  & FER2013    & Hard & Third-party ref    & Off-the-shelf baseline? \\
\bottomrule
\end{tabular}
\end{table}

\textbf{Unified training recipe}: All six experiments share the exact same training configuration (Table 2). The only deliberate exceptions are: (a) VGG16 uses batch size 64 due to VRAM constraints; (b) loss type matches the label type of the training data. This uniformity guarantees that any observed performance difference is caused exclusively by the variable we intentionally changed.

\textbf{Causal chain}:
\begin{enumerate}[label=Step \arabic*:, leftmargin=*]
    \item Exp1 $\to$ Exp1b: Change architecture only (VGG16$\to$ResNet18; same FER2013 data, same recipe).
    \item Exp1b $\to$ Exp2: Change label quality only (FER2013 hard$\to$FERPlus soft; same ResNet18).
    \item Exp2 $\to$ Exp3: Change data diversity only (FERPlus$\to$FERPlus+RAFDB+Self; same ResNet18).
\end{enumerate}

\subsection{Application Deployment}

The best model (ResNet18 + F+R+S, 64.5\% self accuracy) is deployed as EmotionLens, a real-time web application. The backend uses FastAPI + uvicorn + WebSocket: webcam frames stream to the server, YuNet face detector locates faces, ResNet18 infers 7-class probabilities, and EMA temporal smoothing produces stable output. Five independent analysis lenses consume the same emotion stream: (1) crowd emotion statistics (example: cafeteria satisfaction), (2) emotion alert (example: classroom anomaly detection), (3) emotion timeline (example: cinema reaction tracking), (4) nervousness index (example: speech coaching), and (5) mimic game. The frontend uses vanilla HTML/CSS/JS with Swiper.js for lens switching. Public access is provided via Cloudflare Tunnel (HTTPS required for browser \texttt{getUserMedia}).

These five lenses are not the product---they are proofs that the engine generalizes. The engine outputs raw 7-dimensional probability vectors. Anyone can write a new lens that consumes this stream without touching the model. The engine is reusable; the lenses are swappable; the design space is open.

%=====================================================================
\section{Results and Analysis}
%=====================================================================

\subsection{Control-Variable Chain}

Figure \ref{fig:chain} presents the complete control-variable chain. Each step changes exactly one variable; the change in self-data accuracy is the net contribution of that variable.

\begin{figure}[H]
\centering
\includegraphics[width=0.9\textwidth]{control_variable_chain.png}
\caption{Control-variable chain: four-step evolution from course baseline to our best model}
\label{fig:chain}
\end{figure}

\begin{itemize}
    \item \textbf{VGG16 + FER2013}: 11.6\% self --- the course baseline. Effectively non-functional on real faces.
    \item \textbf{ResNet18 + FER2013}: 9.2\% self ($-$2.3pp). Architecture change alone yields no improvement---in fact, a slight regression. Architecture is not the bottleneck.
    \item \textbf{ResNet18 + FERPlus}: 44.4\% self (+35.2pp). Switching from hard to 10-voter soft labels provides the single largest gain. Label quality is the dominant factor.
    \item \textbf{ResNet18 + F+R+S}: 64.5\% self (+20.0pp). Adding real-world data (RAFDB + self-captured) provides the second key gain, recovering negative-emotion recognition.
\end{itemize}

Total improvement: +52.9pp (11.6\% $\to$ 64.5\%).

\subsection{Architecture Ablation}

Table \ref{tab:arch} and Figure \ref{fig:arch} compare VGG16 and ResNet18 on identical FER2013 data.

\begin{table}[H]
\centering
\caption{Architecture ablation: VGG16 vs ResNet18 (same data, same recipe)}
\label{tab:arch}
\begin{tabular}{lcc}
\toprule
Metric & VGG16 (138M) & ResNet18 (11M) \\
\midrule
FER2013 Test Accuracy & 69.46\% & 69.94\% \\
Self Data Accuracy    & 11.58\% & 9.24\% \\
\bottomrule
\end{tabular}
\end{table}

\begin{figure}[H]
\centering
\includegraphics[width=0.85\textwidth]{architecture_ablation.png}
\caption{Architecture ablation: a more modern architecture does not improve real-face performance}
\label{fig:arch}
\end{figure}

On the standard test set, the difference is negligible (0.48pp). On real faces, ResNet18 performs \textit{worse} than VGG16 by 2.34pp. The intuition that ``a better architecture will improve real-world generalization'' does not hold in this setting.

\subsection{Per-Class Recall Analysis}

Figure \ref{fig:recall} shows per-class recall on self data for all four trained models.

\begin{figure}[H]
\centering
\includegraphics[width=\textwidth]{per_class_recall_all.png}
\caption{Per-class recall on self data across four models}
\label{fig:recall}
\end{figure}

The FERPlus-only model achieves \textbf{exactly 0\% recall} on sadness, disgust, and fear. The model literally cannot recognize these emotions on real faces. After adding RAF-DB and self-captured data (F+R+S), recall on these three classes recovers to 53.1\%, 52.7\%, and 59.5\%, respectively. The value of data diversity is concentrated in negative emotions: the model cannot recognize facial expressions it has never seen examples of; real-world training data bridges this gap.

We observe that even after multiple rounds of improvement, model performance on our self-captured test set remains markedly lower than on the FER-provided test sets. However, a closer look at per-class recall reveals an uneven picture: neutral (78\%) and surprise (70\%) perform reasonably well, while disgust (53\%), anger (56\%), and fear (60\%) lag behind considerably. We hypothesize two contributing factors. First, basic expressions such as neutral and happiness are biologically rooted and expressed with high cross-cultural consistency---a smile is a smile; negative emotions, by contrast, are more abstract and subject to greater individual and cultural variation in both expression and interpretation (Jack et al., 2012). Second, the training data (FERPlus, RAF-DB) were annotated predominantly under Western cultural norms, whereas our self-captured test data come entirely from Asian participants. For culturally mediated expressions like disgust, the mismatch between Western annotation standards and Asian facial expressions introduces an inherent ambiguity into the ground truth itself. In other words, the low recall on these emotion classes may reflect not a failure of model learning, but a structural ceiling on label consistency.

\begin{figure}[H]
\centering
\includegraphics[width=0.75\textwidth]{self_confusion_matrix.png}
\caption{Confusion matrix of the best model (F+R+S) on self data}
\label{fig:cm}
\end{figure}

\subsection{Cross-Domain Evaluation}

Table \ref{tab:cross} and Figure \ref{fig:radar} present a complete cross-domain evaluation across all four trained models and the FER library baseline on three test sets.

\begin{table}[H]
\centering
\caption{Cross-domain evaluation: all models $\times$ all test sets}
\label{tab:cross}
\begin{tabular}{lccc}
\toprule
Model & FERPlus Test & FER2013 Test & Self Data \\
\midrule
VGG16 (FER2013)      & 7.1\%  & 69.5\% & 11.6\% \\
ResNet18 (FER2013)   & 9.0\%  & 69.9\% & 9.2\%  \\
ResNet18 (FERPlus)   & 79.9\% & 8.5\%  & 44.4\% \\
ResNet18 (F+R+S)     & \textbf{89.0\%} & 7.0\%  & \textbf{64.5\%} \\
EfficientNet-B0 (FERPlus) & 41.0\% & 10.6\% & 19.4\% \\
\midrule
FER Library          & ---    & ---    & 52.6\% \\
\bottomrule
\end{tabular}
\end{table}

\begin{figure}[H]
\centering
\includegraphics[width=0.8\textwidth]{cross_domain_radar.png}
\caption{Cross-domain radar chart: only the multi-source model generalizes across domains}
\label{fig:radar}
\end{figure}

Each model performs best on its own training domain---a predictable pattern. The key finding is that only the multi-source model (F+R+S) achieves usable performance on real faces (64.5\%). The FER library (hustvl/fer, also trained on FER2013) scores 52.6\% on self data overall, but with disgust recall at 0.0\% and anger at 1.8\%---a pattern identical to our VGG16+FER2013 baseline. This confirms that the problem lies in the data, not our implementation.

\subsection{Domain Gap}

Figure \ref{fig:gap} quantifies the gap between laboratory and real-face performance for each model.

\begin{figure}[H]
\centering
\includegraphics[width=0.85\textwidth]{domain_gap.png}
\caption{Domain gap: lab accuracy vs. real-face accuracy}
\label{fig:gap}
\end{figure}

The domain gap shrinks from 57.9pp (VGG16+FER2013) to 24.5pp (ResNet18+F+R+S). Adding real-world data reduces the lab-to-real gap by more than half.

\subsection{Contribution Decomposition}

Figure \ref{fig:contrib} decomposes the 52.9pp total improvement into three independent factors.

\begin{figure}[H]
\centering
\includegraphics[width=0.85\textwidth]{contribution_breakdown.png}
\caption{Factor contribution decomposition}
\label{fig:contrib}
\end{figure}

\begin{itemize}
    \item \textbf{Architecture} (VGG16$\to$ResNet18): $-$2.3pp --- no contribution; slightly negative.
    \item \textbf{Label quality} (hard$\to$10-voter soft): +35.2pp --- the single largest lever.
    \item \textbf{Data diversity} (single$\to$multi-source): +20.0pp --- second key factor; recovers negative emotions.
    \item \textbf{Net}: +52.9pp (11.6\%$\to$64.5\%).
\end{itemize}

\subsection{FER Library Baseline}

We evaluate the open-source \texttt{hustvl/fer} library on our self data as a third-party reference. Overall accuracy: 52.6\%. Per-class recall: disgust 0.0\%, anger 1.8\%, sadness 3.6\%, fear 12.1\%. This pattern---neutral and happiness recognized adequately, negative emotions essentially invisible---is identical to our VGG16+FER2013 baseline. The consistency across two independent implementations trained on the same data source reinforces our conclusion: the bottleneck is FER2013's label quality and domain coverage, not any particular architectural or implementation choice.

%=====================================================================
\section{Conclusion}
%=====================================================================

We built EmotionLens, a cross-domain facial expression recognition system whose core is a ResNet18 model trained on three complementary data sources: FERPlus soft labels, RAF-DB real-world faces, and a self-captured dataset of 1,123 manually reviewed frames. Through six strictly controlled experiments---each altering exactly one variable---we improved real-face accuracy from 11.6\% (the FER2013+VGG16 course baseline) to 64.5\%, and decomposed the 52.9pp gain into three independent factors: architecture ($-$2.3pp, essentially zero contribution), label quality (+35.2pp, the single largest lever), and data diversity (+20.0pp, the key to recovering negative-emotion recognition).

The contribution of this work is not the 64.5\% number itself, but the causal decomposition that explains \textit{why}. For anyone building FER systems, the actionable takeaway is clear: if you have one improvement budget, spend it on label quality; if you have two, also invest in data diversity. Architecture upgrades alone are unlikely to move the needle on real-world generalization.

EmotionLens's five analysis lenses---crowd emotion statistics, emotion alert, emotion timeline, speech coach, and mimic game---are demonstrations, not products. Each lens consumes the same inference engine; none required modifying the model. The engine is the reusable core. A teacher could replace the lenses with classroom engagement tracking; a therapist could add a session emotion log. The model is the enabler; the applications are limited only by the user's imagination.

\section*{References}

\begin{enumerate}
    \item Barsoum, E., Zhang, C., Ferrer, C. C., \& Zhang, Z. (2016). Training deep networks for facial expression recognition with crowd-sourced label distribution. \textit{Proceedings of the 18th ACM International Conference on Multimodal Interaction (ICMI)}, 279--283.
    \item Chen, T., Pu, T., Wu, H., Xie, Y., Liu, L., \& Lin, L. (2022). Cross-domain facial expression recognition: A unified evaluation benchmark and adversarial graph learning. \textit{IEEE Transactions on Pattern Analysis and Machine Intelligence}, 44(12), 9887--9903.
    \item Ekman, P., \& Friesen, W. V. (1971). Constants across cultures in the face and emotion. \textit{Journal of Personality and Social Psychology}, 17(2), 124--129.
    \item Goodfellow, I. J., Erhan, D., Carrier, P. L., Courville, A., Mirza, M., Hamner, B., ... \& Bengio, Y. (2013). Challenges in representation learning: A report on three machine learning contests. \textit{International Conference on Neural Information Processing (ICONIP)}, 117--124.
    \item He, K., Zhang, X., Ren, S., \& Sun, J. (2016). Deep residual learning for image recognition. \textit{Proceedings of the IEEE Conference on Computer Vision and Pattern Recognition (CVPR)}, 770--778.
    \item Jack, R. E., Garrod, O. G. B., Yu, H., Caldara, R., \& Schyns, P. G. (2012). Facial expressions of emotion are not culturally universal. \textit{Proceedings of the National Academy of Sciences (PNAS)}, 109(19), 7241--7244.
    \item Li, S., \& Deng, W. (2019). Reliable crowdsourcing and deep locality-preserving learning for unconstrained facial expression recognition. \textit{IEEE Transactions on Image Processing}, 28(1), 356--370.
    \item Loshchilov, I., \& Hutter, F. (2019). Decoupled weight decay regularization. \textit{International Conference on Learning Representations (ICLR)}.
    \item Simonyan, K., \& Zisserman, A. (2014). Very deep convolutional networks for large-scale image recognition. \textit{arXiv preprint arXiv:1409.1556}.
    \item Szegedy, C., Vanhoucke, V., Ioffe, S., Shlens, J., \& Wojna, Z. (2016). Rethinking the inception architecture for computer vision. \textit{Proceedings of the IEEE Conference on Computer Vision and Pattern Recognition (CVPR)}, 2818--2826.
    \item Tan, M., \& Le, Q. V. (2019). EfficientNet: Rethinking model scaling for convolutional neural networks. \textit{Proceedings of the 36th International Conference on Machine Learning (ICML)}, 6105--6114.
    \item Zhang, H., Cisse, M., Dauphin, Y. N., \& Lopez-Paz, D. (2018). mixup: Beyond empirical risk minimization. \textit{International Conference on Learning Representations (ICLR)}.
    \item Zhang, K., Zhang, Z., Li, Z., \& Qiao, Y. (2016). Joint face detection and alignment using multitask cascaded convolutional networks. \textit{IEEE Signal Processing Letters}, 23(10), 1499--1503.
\end{enumerate}

\end{document}
